\documentclass[11pt]{article}
\usepackage[margin=1in]{geometry}
\usepackage{amsmath,booktabs}
\usepackage[hidelinks]{hyperref}
\title{Faster Convergence of Sparse Attention\\through Adaptive Token Pruning}
\author{Ada Lindqvist \and Rafael Moreno}
\date{}
\begin{document}
\maketitle
\begin{abstract}
Transformer models spend most of their compute on attention over tokens that contribute little to the
final prediction. We propose adaptive token pruning, a simple method that removes low-importance tokens
between layers. On three language modelling benchmarks our method reduces training time by 31\% while
matching the perplexity of a dense baseline.
\end{abstract}

\section{Introduction}
Large language models are expensive to train, and a large share of that cost comes from the attention
mechanism, whose complexity grows quadratically with sequence length. Many tokens, such as punctuation
and function words, receive little attention in deeper layers, yet they are processed at every step.

Earlier work has explored fixed sparsity patterns and learned routing. Fixed patterns are cheap but
ignore the content of the sequence, while learned routing adds parameters and is hard to train. We take
a middle path: we score every token with a lightweight importance estimate and drop the lowest-scoring
tokens after each block.

Our contributions are a pruning rule that needs no extra parameters, an analysis of its effect on
gradient flow, and experiments showing faster convergence on three benchmarks.

\section{Method}
For a sequence of hidden states $h_1, \dots, h_n$ at layer $\ell$, we define the importance of token $i$
as the total attention it receives:
\begin{equation}
  s_i^{(\ell)} = \sum_{j=1}^{n} A^{(\ell)}_{ji}.
\end{equation}
After each block we keep the $k = \lceil \rho n \rceil$ tokens with the highest scores and pass only those
to the next layer. The keep ratio $\rho$ is fixed at 0.7 for all layers. Pruned tokens are restored at the
output layer by copying their last hidden state, so the model still produces a prediction for every position.

\section{Experiments}
We train 350M-parameter models on WikiText-103, The Pile and C4 for 50{,}000 steps with a batch size of
256 sequences. Table~\ref{tab:main} compares perplexity and wall-clock training time.

\begin{table}[h]
\centering
\caption{Perplexity and training time on three benchmarks.}
\label{tab:main}
\begin{tabular}{lccc}
\toprule
Model & WikiText-103 & The Pile & C4 \\
\midrule
Dense baseline & 18.4 & 12.9 & 16.2 \\
Adaptive pruning (ours) & 18.6 & 13.1 & 16.3 \\
\midrule
Training time (hours) & 41 / 29 & 96 / 66 & 88 / 61 \\
\bottomrule
\end{tabular}
\end{table}

Pruning costs less than 0.2 perplexity points on every benchmark while cutting training time by about
a third. The gap is largest on The Pile, which contains long documents with many repeated tokens.

\section{Conclusion}
Adaptive token pruning is a simple way to make attention cheaper without new parameters. In future work
we plan to apply it to encoder models.
\end{document}
